\section{Empirical Validation of Foundational Assumptions}
\label{sec:empirical_validation_assumptions}

The analysis in this chapter is built upon a set of foundational assumptions
regarding the distribution of grammatical and semantic information across tokens,
sequences, and training batches.
While these assumptions are motivated by linguistic structure and common training
practice, their validity is ultimately an empirical question.
The purpose of the experiments in this section is therefore to quantitatively
evaluate these assumptions in realistic transformer models and datasets.

Specifically, we aim to measure (i) the degree to which grammatical structure induces
globally shared, low-dimensional patterns in activations and gradients; (ii) the
strength of semantic correlation among tokens within the same sequence; and (iii)
the extent to which semantic correlations diminish across independently shuffled
sequences or sufficiently distant positions in long-context settings.
By explicitly characterizing these correlation patterns, we seek to verify that
simple averaging across appropriately chosen tokens suppresses independent semantic
variation while preserving the shared grammatical structure that dominates gradient
geometry.
This empirical validation is critical for establishing when and why the proposed
averaging-based approximation accurately captures the spiky, low-rank structure
driving optimization dynamics.

\paragraph{Preliminary spectral alignment observation.}
Given an activation tensor $a\in\mathbb{R}^{\texttt{B}\times\texttt{L}\times\texttt{H}}$, we reshape it into
$b\in\mathbb{R}^{(\texttt{B}\cdot\texttt{L})\times\texttt{H}}$ and form a ``column-mean matrix''
$\mathrm{ColMean}(b)$ by averaging each column of $b$ and broadcasting the resulting mean vector to match the shape of $b$.
We then compute the SVD of $b$ and measure the (absolute) cosine similarity between $\mathrm{ColMean}(b)$ and the reconstruction induced by the $k$-th singular component.
Empirically, we observe that the leading $\approx 1.5\%$ of the spectral subspace exhibits high alignment with $\mathrm{ColMean}(b)$,
whereas the remaining $\approx 98.5\%$ is nearly orthogonal.
Moreover, the alignment strength depends on how the column mean is estimated: averaging over all sentences in the batch yields the
strongest alignment, followed by averaging over randomly sampled tokens across sentences, while using a single sentence produces the
weakest alignment.

\paragraph{Spectral alignment under different mean estimators.}
We formalize and extend the above observation by comparing three estimators of the column mean.
Let $a\in\mathbb{R}^{B\times L\times H}$ denote the activation tensor, and let $b\in\mathbb{R}^{(BL)\times H}$ be the reshaped matrix obtained by
flattening the batch and sequence dimensions.
Given a mean vector $\mu\in\mathbb{R}^{H}$, we define the corresponding broadcasted ``column-mean matrix'' as
$\mathrm{ColMean}_{\mu}(b)=\mathbf{1}\mu^{\top}\in\mathbb{R}^{(BL)\times H}$.
For each $k$, we compare $\mathrm{ColMean}_{\mu}(b)$ with the rank-$1$ reconstruction induced by the $k$-th singular component of $b$, i.e.,
$\sigma_k u_k v_k^{\top}$, using the absolute cosine similarity in Frobenius inner product.

We consider the following three choices of $\mu$:
(i) \emph{Batch mean} (Method~1): average each column of $b$ over all $BL$ tokens in the full batch;
(ii) \emph{Single-sequence mean} (Method~2): compute $\mu$ from one sequence only (e.g., averaging columns of $a[0,:,:]\in\mathbb{R}^{L\times H}$) and broadcast it to shape $(BL)\times H$;
(iii) \emph{Random token sampling across sequences} (Method~3): randomly sample a subset of tokens from each of the $B$ sequences to form a
set of token vectors, compute the column mean over this set, and broadcast.

Across all three estimators, we find a consistent spectral dichotomy: the leading $\approx 1.5\%$ of singular components exhibits high
alignment with $\mathrm{ColMean}_{\mu}(b)$, while the remaining $\approx 98.5\%$ is nearly orthogonal.
Moreover, the alignment strength follows a clear ordering,
\emph{batch mean} (Method~1) $>$ \emph{random token sampling} (Method~3) $>$ \emph{single-sequence mean} (Method~2),
which is consistent with the chapter's core hypothesis that the shared, low-rank structure is most reliably recovered when averaging over
many semantically diverse tokens while preserving the globally shared grammatical signal.
Figures~\ref{fig:colmean_alignment_method1}--\ref{fig:colmean_alignment_method3} report the cosine similarity curves for the three estimators.

\begin{figure}[t]
  \centering
  \includegraphics[width=0.92\linewidth]{figures/col_batch.png}
  \caption{Absolute cosine similarity between $\mathrm{ColMean}_{\mu}(b)$ (with $\mu$ estimated using the full batch) and the rank-$1$ matrix
  $\sigma_k u_k v_k^{\top}$ induced by the $k$-th singular component of $b$.}
  \label{fig:colmean_alignment_method1}
\end{figure}

\begin{figure}[t]
  \centering
  \includegraphics[width=0.92\linewidth]{figures/col_single.png}
  \caption{Same as Figure~\ref{fig:colmean_alignment_method1}, but with $\mu$ estimated from a single sequence (e.g., $a[0,:,:]$) and then
  broadcast to shape $(BL)\times H$.}
  \label{fig:colmean_alignment_method2}
\end{figure}

\begin{figure}[t]
  \centering
  \includegraphics[width=0.92\linewidth]{figures/col_sampled.png}
  \caption{Same as Figure~\ref{fig:colmean_alignment_method1}, but with $\mu$ estimated from randomly sampled tokens across sequences.}
  \label{fig:colmean_alignment_method3}
\end{figure}